\section{Developmental experiments}
\label{app:developmental}

Experiments~1--3 developed the modeling and computational workflow but do not
contribute numerical evidence to the main comparison.

\begin{table}[htbp]
\centering
\small
\caption{Developmental progression to Experiment~4.}
\begin{tabularx}{\textwidth}{
  >{\raggedright\arraybackslash}p{0.12\textwidth}
  >{\raggedright\arraybackslash}p{0.25\textwidth}
  >{\raggedright\arraybackslash}p{0.27\textwidth}
  >{\raggedright\arraybackslash}X}
\toprule
Experiment & Purpose & Why developmental & Link to Experiment~4 \\
\midrule
1
& Construct and illustrate the Gamma-noise model under constant and
piecewise transmission scenarios.
& Each scenario used one simulated epidemic and primarily examined model
construction and starting-value sensitivity.
& Established the basic latent-rate recovery workflow. \\
2
& Exercise a high-particle, multi-start Gamma-noise fitting workflow.
& The workflow was applied to one fixed simulated data set.
& Established the repeated-fitting workflow. \\
3
& Evaluate Gamma-noise recovery across 200 tasks with
\(N_{\mathrm{mif}}=100\).
& Its lower-iteration analysis was superseded for final numerical claims.
& Motivated the \(N_{\mathrm{mif}}=600\) paired comparison. \\
4
& Compare Gamma-noise and constant-\(B\) recovery on shared data.
& Main experiment rather than developmental evidence.
& Supplies all numerical results in the main text. \\
\bottomrule
\end{tabularx}
\end{table}
